\documentclass[11pt,a4paper]{article}
\usepackage{oi-statement}
\OISetup{CSP-J 模拟赛}{第四套}{入门级}
\OISetCompileOptions{-O2 -std=c++14}
\begin{document}
\OICoverTitle
\OICoverTable{4}{
题目名称 & 接龙 & 平方和 & 彩彩的三彩项链 & 表达式\\\hline

题目类型 & 传统型 & 传统型 & 传统型 & 传统型\\\hline

目录 & {\ttfamily chain} & {\ttfamily sumsq} & {\ttfamily necklace} & {\ttfamily eval}\\\hline

可执行文件名 & {\ttfamily chain} & {\ttfamily sumsq} & {\ttfamily necklace} & {\ttfamily eval}\\\hline

输入文件名 & {\ttfamily chain.in} & {\ttfamily sumsq.in} & {\ttfamily necklace.in} & {\ttfamily eval.in}\\\hline

输出文件名 & {\ttfamily chain.out} & {\ttfamily sumsq.out} & {\ttfamily necklace.out} & {\ttfamily eval.out}\\\hline

每个测试点时限 & 1.0 秒 & 1.0 秒 & 1.0 秒 & 1.0 秒\\\hline

内存限制 & 512 MiB & 512 MiB & 512 MiB & 512 MiB\\\hline

测试点数目 & 10 & 10 & 10 & 10\\\hline

测试点是否等分 & 是，每点 10 分 & 是，每点 10 分 & 是，每点 10 分 & 是，每点 10 分\\\hline
}
\OICompilation{4}{
对于 C++ 语言 & {\ttfamily chain.cpp} & {\ttfamily sumsq.cpp} & {\ttfamily necklace.cpp} & {\ttfamily eval.cpp}\\\hline
}
\OICoverNotes
\begin{enumerate}
\item 本套为 CSP-J 训练模拟赛，共 4 题，每题 100 分，总分 400 分。
\item 各题均采用标准输入输出，提交对应的同名 C++ 源程序。
\item 每题有 10 个测试点，每点 10 分，按测试点独立计分。
\item 答案采用普通文本比较，忽略行末空格和文件末尾多余换行。
\item 题面样例与下发大样例不计入正式测试点。
\item 程序可使用的栈空间与题目的内存限制一致。
\item 样例框左侧行号不属于输入输出。
\item 使用 C++14，编译选项为 \texttt{-O2 -std=c++14}。
\end{enumerate}

\OIProblem{接龙}{chain}

\OISection{题目描述}

牛牛有 $n$ 张牌，第 $i$ 张牌上的数字为 $a_i$。他可以任意调整这些牌的放置顺序，然后依次把牌放到桌面一列的末尾。

每放置一张牌后，他可以进行至多一次收牌操作：从桌面上选择两张数字相同的牌，收走这两张牌以及它们之间的所有牌。若这两张相同的牌数字为 $x$，本次共收走 $len$ 张牌，则获得

\[
x\times len
\]

分。被收走的牌不会再次使用。牛牛也可以在某一回合不收牌。

求把所有牌都放置完以后，能够获得的最大总分。

\OISection{输入格式}

第一行输入一个正整数 $n$。

第二行输入 $n$ 个正整数 $a_1,a_2,\ldots,a_n$。

\OISection{输出格式}

输出一行一个整数，表示最大总分。

\begin{OISample}{1}{输入}
4
3 3 4 4
\end{OISample}
\begin{OISample}{1}{输出}
16
\end{OISample}
\begin{OISample}{2}{输入}
3
1 2 3
\end{OISample}
\begin{OISample}{2}{输出}
0
\end{OISample}
\begin{OISample}{3}{输入}
3
5 5 5
\end{OISample}
\begin{OISample}{3}{输出}
15
\end{OISample}
\OIAdditionalSample{4}{chain/chain}
\OIAdditionalSample{5}{chain/chain}
\Needspace{9\baselineskip}
\OISection{数据范围}

\begin{itemize}
\item $1\le n\le10^5$；
\item $1\le a_i\le10^5$。
\end{itemize}

本题共 $10$ 个测试点，每点 $10$ 分，下列前几类独立计分。

\begin{OIDataTable}{|c|c|c|}
测试点 & 特殊限制 & 分值\\\hline
$1\sim2$ & $n\le100$ & $20$ 分\\\hline
$3\sim4$ & 所有牌数字相同 & $20$ 分\\\hline
$5\sim7$ & $a_i\le1000$ & $30$ 分\\\hline
$8\sim10$ & 无额外限制 & $30$ 分\\\hline
\end{OIDataTable}

\OIProblem{平方和}{sumsq}

\OISection{题目描述}

给定 $n$ 个正整数。请选择两个\textbf{不同位置}的数字 $a_i,a_j$，令

\[
x=a_i^2+a_j^2\text{。}
\]

要求 $x$ 能被 $k$ 整除，并使 $x$ 尽可能大。输入保证至少存在一组合法选择。

\OISection{输入格式}

第一行包含两个正整数 $n,k$。

第二行包含 $n$ 个正整数 $a_1,a_2,\ldots,a_n$。

\OISection{输出格式}

输出最大的合法平方和。

\begin{OISample}{1}{输入}
5 2
1 2 3 4 5
\end{OISample}
\begin{OISample}{1}{输出}
34
\end{OISample}
\begin{OISample}{2}{输入}
5 10
1 2 3 4 5
\end{OISample}
\begin{OISample}{2}{输出}
20
\end{OISample}
\OIAdditionalSample{3}{sumsq/sumsq}
\OIAdditionalSample{4}{sumsq/sumsq}
\Needspace{10\baselineskip}
\OISection{数据范围}

\begin{itemize}
\item $1\le n\le2\times10^5$；
\item $1\le k\le10^9$；
\item $1\le a_i\le10^9$；
\item 保证存在两个不同位置组成合法答案。
\end{itemize}

本题共 $10$ 个测试点，每点 $10$ 分，下列前几类独立计分。

\begin{OIDataTable}{|c|c|c|}
测试点 & 特殊限制 & 分值\\\hline
$1\sim2$ & $n\le2000$ & $20$ 分\\\hline
$3\sim5$ & $k\le10^5$ & $30$ 分\\\hline
$6\sim10$ & 无额外限制 & $50$ 分\\\hline
\end{OIDataTable}

\OIProblem{彩彩的三彩项链}{necklace}

\OISection{题目描述}

彩彩有一条首尾相连、长度为 $n$ 的项链。每颗珠子的颜色是红、绿、蓝之一，分别用字符 \texttt{\detokenize{r}}、\texttt{\detokenize{g}}、\texttt{\detokenize{b}} 表示。

一次点击可以把一颗珠子的颜色按照以下顺序改变一次：

\begin{quote}\ttfamily
r -> g -> b -> r
\end{quote}

可以多次点击同一颗珠子。

彩彩希望修改后，项链上任意两颗相邻珠子的颜色都不同。由于项链首尾相连，第 $1$ 颗珠子和第 $n$ 颗珠子也相邻。

求最少点击次数。

\OISection{输入格式}

第一行输入一个整数 $n$。

第二行输入一个长度为 $n$ 的字符串 $s$。

\OISection{输出格式}

输出最少点击次数。

\begin{OISample}{1}{输入}
3
rrr
\end{OISample}
\begin{OISample}{1}{输出}
3
\end{OISample}
\OIAdditionalSample{2}{necklace/necklace}
\OIAdditionalSample{3}{necklace/necklace}
\Needspace{10\baselineskip}
\OISection{数据范围}

\begin{itemize}
\item $2\le n\le10^6$；
\item $s$ 只包含 \texttt{\detokenize{r}}、\texttt{\detokenize{g}}、\texttt{\detokenize{b}}。
\end{itemize}

本题共 $20$ 个测试点，每点 $5$ 分，下列各组独立计分。

\begin{OIDataTable}{|c|c|c|}
测试点 & 特殊性质 & 分值\\\hline
$1\sim2$ & $n=10$ & $10$ 分\\\hline
$3\sim4$ & $n=1000$ & $10$ 分\\\hline
$5\sim6$ & 字符串只含一种颜色 & $10$ 分\\\hline
$7\sim8$ & 字符串只含两种颜色 & $10$ 分\\\hline
$9\sim20$ & 无额外限制 & $60$ 分\\\hline
\end{OIDataTable}

\OIProblem{表达式}{eval}

\OISection{题目描述}

给定一个只由正整数、加号 \texttt{\detokenize{+}} 和乘号 \texttt{\detokenize{*}} 组成的合法表达式 $S$。

你可以在表达式中的任意合法位置添加任意多对括号，但不能改变数字和运算符的原有顺序。求能够得到的最大表达式值。

\OISection{输入格式}

一行一个字符串 $S$。

\OISection{输出格式}

输出一行一个正整数，表示最大值。答案可能非常大，需要输出完整十进制表示。

\begin{OISample}{1}{输入}
1*2+2*1
\end{OISample}
\begin{OISample}{1}{输出}
4
\end{OISample}
\begin{OISample}{2}{输入}
1*2345+6789*9867+5432*1
\end{OISample}
\begin{OISample}{2}{输出}
139741066
\end{OISample}
\OISection{样例说明}

样例 1 可以变为 $1\times(2+2)\times1=4$。

样例 2 可以变为

\[
(1\times2345+6789)\times(9867+5432\times1)
=1\times(2345+6789)\times(9867+5432)\times1\text{。}
\]

\OIAdditionalSample{3}{eval/eval}
\OIAdditionalSample{4}{eval/eval}
\Needspace{15\baselineskip}
\OISection{数据范围}

设 $|S|$ 为字符串长度，$n$ 为正整数个数，第 $i$ 个数为 $a_i$。

\begin{itemize}
\item $1\le|S|\le5\times10^3$；
\item $1\le n\le\lceil |S|/2\rceil$；
\item $1\le a_i<10^{|S|}$。
\end{itemize}

本题共 $20$ 个测试点，每点 $5$ 分。

\begin{OIDataTable}{|c|c|c|c|}
测试点 & $|S|\le$ & $a_i$ & 特殊性质\\\hline
$1\sim2$ & $10$ & $<10^9$ & 无\\\hline
$3\sim4$ & $5000$ & $\le1$ & 无\\\hline
$5\sim6$ & $5000$ & $\le9$ & 无\\\hline
$7\sim8$ & $5000$ & $<10^9$ & 只有 \texttt{\detokenize{+}}\\\hline
$9\sim10$ & $5000$ & $<10^{|S|}$ & 只有 \texttt{\detokenize{+}}\\\hline
$11\sim12$ & $5000$ & $<10^{|S|}$ & 只有 \texttt{\detokenize{*}}\\\hline
$13\sim14$ & $1000$ & $<10^{|S|}$ & 无\\\hline
$15\sim20$ & $5000$ & $<10^{|S|}$ & 无\\\hline
\end{OIDataTable}

\end{document}
